% Generated draft: requires mathematical and source review.
\documentclass[11pt]{article}
\usepackage[margin=25mm]{geometry}
\usepackage{fontspec}
\setmainfont{Latin Modern Roman}
\usepackage{amsmath,amssymb,mathtools,unicode-math}
\setmathfont{Latin Modern Math}
\usepackage{enumitem,needspace,xurl,hyperref}
\setlength{\parindent}{0pt}
\setlength{\parskip}{5pt}
\setlength{\emergencystretch}{3em}
\newcommand{\N}{\mathbb{N}}
\newcommand{\Z}{\mathbb{Z}}
\newcommand{\Q}{\mathbb{Q}}
\newcommand{\R}{\mathbb{R}}
\newcommand{\C}{\mathbb{C}}
\newcommand{\F}{\mathbb{F}}
\newcommand{\A}{\mathbb{A}}
\newcommand{\PP}{\mathbb{P}}
\newcommand{\E}{\mathbb{E}}
\newcommand{\eps}{\varepsilon}
\newcommand{\dd}{\,\mathrm{d}}
\begin{document}
\section*{UnsolvedMath ID: 30002910}
\textbf{Source problem code:} catalogue entry
\subsection{Definitions and mathematical statement}
A hyperbolic \(n\)-manifold is a Riemannian \(n\)-manifold of constant sectional curvature \(-1\). Completeness refers to its Riemannian metric. A finite-volume cusp is an end with a compact flat horospherical cross-section; write \(\kappa(M)\) for the number of such cusps and \(\operatorname{Vol}(M)\) for hyperbolic volume.

For an integer dimension \(n\), a positive integer \(c\), and a finite real number \(V>0\), the source asks about the counting function formally written as
\[
\rho_n^c(V)=
\#\left(
\left\{M:\ M\text{ is a complete hyperbolic }n\text{-manifold},\quad
\kappa(M)=c,\quad \operatorname{Vol}(M)<V\right\}/\!\sim
\right).
\]
Here \(\sim\) denotes the equivalence relation under which manifolds are counted. The supplied record does not specify this relation; counting up to isometry is a possible interpretation, but is not asserted as a source hypothesis. The volume inequality is strict and implies finite volume. The source also leaves the intended dimension range, connectedness, orientability, and treatment of boundary unspecified. These conventions must be resolved before this expression defines a unique counting problem.

\textbf{Source question.} What can be established about \(\rho_n^c(V)\), for positive fixed cusp number \(c\), as the dimension and volume bound vary? In particular, the supplied context concerns its growth as \(V\to\infty\). The source gives a research question rather than a specific conjectured bound, asymptotic formula, or quantified theorem to prove.

\textbf{Supplied qualification.} The dated curation labels the problem partially solved and reports four-dimensional fixed-cusp growth of the form \(C^{V\log V}\), while reporting that a complete general-dimensional description remains unavailable. This is a supplied literature assessment, not an independently verified mathematical assertion. The record does not specify the meaning of this growth comparison, the constants and their dependence, or the cusp numbers for which it is asserted.
\subsection{Short English statement}
For a fixed positive cusp number \(c\), study how many complete hyperbolic \(n\)-manifolds have exactly \(c\) cusps and volume strictly below \(V\), especially their growth as \(V\) increases. The counting equivalence and several manifold conventions remain unspecified. The supplied dated assessment reports partial progress in dimension four and requires independent review.
\subsection{Sources}
\begin{itemize}
\item Supplied local UnsolvedMath registry record, permanent ID \(30002910\), title \emph{Counting Hyperbolic Manifolds with Fixed Cusp Number}, exact source problem code \texttt{catalogue entry}, in the Oberwolfach Reports Open Problems collection. No individual proposer is identified; the proposed year is recorded as 2015.
\item Supplied report citation: \emph{New Perspectives on the Interplay between Discrete Groups in Low-Dimensional Topology and Arithmetic Lattices} (2016), Oberwolfach Reports, pp.~1701--1746, DOI \texttt{10.4171/owr/2015/30}. Original report link: \url{https://doi.org/10.4171/owr/2015/30}. The background describes Oberwolfach Reports as documenting mathematical workshops and their open-problem discussions.
\item Supplied literature source labeled ``Kolpakov and Martelli primary paper'': \url{https://doi.org/10.1007/s00039-013-0247-2}. The record supplies no paper title, theorem number, or precise counting statement from this reference.
\item Literature assessment checked at \texttt{2026-08-21}: supplied dated curation, not independently verified. It labels the problem partially solved, reports four-dimensional fixed-cusp growth of the form \(C^{V\log V}\), and reports that a complete general-dimensional description remains unavailable. The record explicitly requests independent verification of time-sensitive status and cited claims.
\item The supplied statement status is \texttt{corrected\_verified}; its accompanying verification note says that only safe notation, typography, and TeX-delimiter cleanup was applied, preserving hypotheses, quantifiers, and subquestions. This supplied label does not constitute independent mathematical verification of this draft.
\end{itemize}
\end{document}
